\section{Implementation}

\label{sec:implementation}
% ══════════════════════════════════════════════════════════════════════════════

\subsection{Software Architecture}

ZooMRV is implemented as a set of containerized microservices:

\begin{itemize}
  \item \textbf{Data ingestion}: Google Earth Engine and Planetary
    Computer APIs for satellite data access; custom GEDI processor.
  \item \textbf{CarbonNet inference}: PyTorch 2.2 model serving on
    NVIDIA A100 GPUs; batch processing of Sentinel-2 tiles.
  \item \textbf{Change detection}: Temporal analysis pipeline using
    NumPy, Rasterio, and scikit-image.
  \item \textbf{Baseline engine}: Synthetic control computation using
    SciPy optimization; control area matching with H3 spatial indexing.
  \item \textbf{Smart contracts}: Solidity contracts on Zoo Network;
    IPFS for data storage.
  \item \textbf{API layer}: FastAPI REST API exposing project status,
    credit issuance, and verification endpoints.
\end{itemize}

\subsection{Computational Requirements}

Processing a 1-million-hectare project requires approximately:
\begin{itemize}
  \item 2.4 GPU-hours for CarbonNet inference per monitoring period.
  \item 0.8 CPU-hours for change detection per monitoring period.
  \item 0.2 CPU-hours for baseline computation (one-time setup).
  \item 12\,GB storage per monitoring period (compressed imagery and
    derived products).
\end{itemize}

Annual monitoring of the full 2.3 Mha portfolio costs approximately
\$18{,}000 in compute, or \$0.008 per hectare---orders of magnitude
cheaper than field-based monitoring.


% ══════════════════════════════════════════════════════════════════════════════
